\documentclass[12pt,a4paper]{amsart}

\usepackage{amsmath,amssymb,amsthm}
\usepackage{mathtools}
\usepackage{hyperref}
\usepackage{geometry}
\geometry{margin=2.5cm}

% -------------------------------------------------------
% Theorem environments
% -------------------------------------------------------
\newtheorem{theorem}{Theorem}[section]
\newtheorem{lemma}[theorem]{Lemma}
\newtheorem{corollary}[theorem]{Corollary}
\newtheorem{proposition}[theorem]{Proposition}
\theoremstyle{definition}
\newtheorem{definition}[theorem]{Definition}
\newtheorem{remark}[theorem]{Remark}

% -------------------------------------------------------
% Custom commands
% -------------------------------------------------------
\newcommand{\N}{\mathbb{N}}
\newcommand{\Z}{\mathbb{Z}}
\newcommand{\Q}{\mathbb{Q}}
\newcommand{\R}{\mathbb{R}}
\newcommand{\C}{\mathbb{C}}
\DeclareMathOperator{\Re}{Re}
\DeclareMathOperator{\Im}{Im}
\DeclareMathOperator{\res}{Res}

% -------------------------------------------------------
\begin{document}
% -------------------------------------------------------

\title{The Riemann Zeta Function:\\
       Definition, Analytic Continuation, and the Functional Equation}

\author{Michael Fuhrmann}
\address{Specialist Mathematics Project, Build 84}
\email{specialist-maths@localhost}
\date{\today}

\begin{abstract}
The Riemann zeta function $\zeta(s)$ is one of the most profound objects in
all of mathematics, connecting number theory, complex analysis, and physics.
We develop the theory from first principles: the Dirichlet series definition
valid for $\Re(s) > 1$, the Euler product over primes, the analytic
continuation to $\C \setminus \{1\}$, and Riemann's functional equation
\[
  \zeta(s) \;=\; 2^s\,\pi^{s-1}\,\sin\!\Bigl(\tfrac{\pi s}{2}\Bigr)\,
  \Gamma(1-s)\,\zeta(1-s).
\]
We identify the trivial zeros at $s = -2, -4, -6, \ldots$ and introduce the
Riemann--Siegel formula $Z(t) = e^{i\theta(t)}\,\zeta\!\bigl(\tfrac{1}{2}+it\bigr)$
for numerically locating zeros on the critical line.
The critical strip $0 < \Re(s) < 1$ is the arena in which all
non-trivial zeros reside; the Riemann Hypothesis (still open) asserts
that all of them satisfy $\Re(s) = \tfrac{1}{2}$.
\end{abstract}

\maketitle
\tableofcontents

% ================================================================
\section{Introduction}
% ================================================================

In his celebrated 1859 memoir \emph{\"Uber die Anzahl der Primzahlen unter
einer gegebenen Gr\"o\ss e} \cite{Riemann1859}, Bernhard Riemann introduced
the complex-variable study of the function
\[
  \zeta(s) = \sum_{n=1}^{\infty} \frac{1}{n^s},
\]
which had been considered by Euler only for real $s > 1$.
Riemann's key insight was to extend $\zeta$ to the whole complex plane,
derive its functional equation, and relate its zeros to the distribution
of prime numbers.  This paper develops that foundational theory carefully.

\begin{remark}
We follow the standard convention: $s = \sigma + it$ denotes a complex
number with real part $\sigma = \Re(s)$ and imaginary part $t = \Im(s)$.
The function $\Gamma$ denotes Euler's gamma function, defined in
Section~\ref{sec:gamma} below.
\end{remark}

% ================================================================
\section{The Dirichlet Series for \texorpdfstring{$\Re(s)>1$}{Re(s)>1}}
% ================================================================

\begin{definition}[Riemann Zeta Function --- initial domain]
\label{def:zeta_series}
For $s \in \C$ with $\Re(s) > 1$, the \emph{Riemann zeta function} is
\[
  \zeta(s) \;=\; \sum_{n=1}^{\infty} \frac{1}{n^s}.
\]
\end{definition}

\begin{proposition}[Absolute convergence]
\label{prop:convergence}
The series in Definition~\ref{def:zeta_series} converges absolutely and
uniformly on every compact subset of the half-plane $\Re(s) > 1$.
In particular, $\zeta(s)$ is holomorphic on $\{s : \Re(s) > 1\}$.
\end{proposition}

\begin{proof}
Fix $s = \sigma + it$ with $\sigma > 1$.  Then $|n^{-s}| = n^{-\sigma}$,
and $\sum_{n=1}^\infty n^{-\sigma}$ converges by the integral test
(the integral $\int_1^\infty x^{-\sigma}\,dx = (\sigma-1)^{-1}$ is finite
for $\sigma > 1$).  Uniform convergence on compact subsets follows from
the Weierstrass $M$-test with $M_n = n^{-\sigma_0}$ where
$\sigma_0 = \min_K \Re(s) > 1$.  Holomorphicity then follows from
uniform convergence of the series of holomorphic functions $s \mapsto n^{-s}$.
\end{proof}

% ================================================================
\section{The Euler Product}
% ================================================================

\begin{theorem}[Euler Product]
\label{thm:euler_product}
For $\Re(s) > 1$,
\[
  \zeta(s) \;=\; \prod_{p\;\text{prime}} \frac{1}{1 - p^{-s}},
\]
where the product converges absolutely.
\end{theorem}

\begin{proof}
We use the geometric series: for each prime $p$ and $\Re(s) > 1$,
\[
  \frac{1}{1 - p^{-s}} = \sum_{k=0}^{\infty} p^{-ks}.
\]
Expanding the finite product over primes $p \le P$:
\[
  \prod_{p \le P} \frac{1}{1 - p^{-s}}
  = \prod_{p \le P} \sum_{k=0}^{\infty} p^{-ks}
  = \sum_{\substack{n \ge 1 \\ p \mid n \Rightarrow p \le P}} n^{-s},
\]
where the last equality uses the Fundamental Theorem of Arithmetic
(unique factorization) to identify the terms with integers whose prime
factors are all at most $P$.

The difference between $\zeta(s)$ and this partial product is
\[
  \Bigl|\zeta(s) - \prod_{p \le P}\frac{1}{1-p^{-s}}\Bigr|
  \;\le\; \sum_{n > P} n^{-\sigma},
\]
which tends to $0$ as $P \to \infty$ (since $\sum n^{-\sigma}$ converges
for $\sigma > 1$).  Thus the product over all primes converges to $\zeta(s)$.

Absolute convergence of the Euler product follows from
$\sum_p p^{-\sigma} < \sum_n n^{-\sigma} < \infty$ for $\sigma > 1$.
\end{proof}

\begin{corollary}
$\zeta(s) \ne 0$ for $\Re(s) > 1$.
\end{corollary}

\begin{proof}
An absolutely convergent product of nonzero factors is nonzero.
Each factor $(1 - p^{-s})^{-1} \ne 0$ since $|p^{-s}| = p^{-\sigma} < 1$.
\end{proof}

% ================================================================
\section{The Gamma Function}
\label{sec:gamma}
% ================================================================

\begin{definition}[Gamma function]
\label{def:gamma}
For $\Re(s) > 0$, Euler's \emph{gamma function} is
\[
  \Gamma(s) = \int_0^{\infty} t^{s-1} e^{-t}\,dt.
\]
\end{definition}

\begin{proposition}[Key properties of $\Gamma$]
\label{prop:gamma_props}
\begin{enumerate}
  \item[\textup{(i)}] \textit{Functional equation:} $\Gamma(s+1) = s\,\Gamma(s)$.
  \item[\textup{(ii)}] \textit{Factorial:} $\Gamma(n) = (n-1)!$ for $n \in \N$.
  \item[\textup{(iii)}] \textit{Reflection formula:}
    $\Gamma(s)\,\Gamma(1-s) = \dfrac{\pi}{\sin(\pi s)}$ for $s \notin \Z$.
  \item[\textup{(iv)}] \textit{Meromorphic continuation:} $\Gamma$ extends to a
    meromorphic function on $\C$ with simple poles at $s = 0, -1, -2, \ldots$
\end{enumerate}
\end{proposition}

\begin{proof}[Proof sketch]
(i) Integrate by parts: $\int_0^\infty t^s e^{-t}\,dt = [-t^s e^{-t}]_0^\infty
+ s\int_0^\infty t^{s-1}e^{-t}\,dt = s\,\Gamma(s)$.
(ii) Follows from (i) and $\Gamma(1) = 1$.
(iii) The reflection formula is a classical identity proved via the infinite
product for $\Gamma$ and the Weierstrass product for $\sin$; see \cite{Ahlfors1979}.
(iv) The recursion $\Gamma(s) = \Gamma(s+n)/[s(s+1)\cdots(s+n-1)]$ provides
the meromorphic continuation.
\end{proof}

% ================================================================
\section{Analytic Continuation of \texorpdfstring{$\zeta$}{zeta}}
% ================================================================

\begin{theorem}[Analytic continuation via the Gamma integral]
\label{thm:continuation}
For $\Re(s) > 1$,
\[
  \Gamma(s)\,\zeta(s) = \int_0^{\infty} \frac{t^{s-1}}{e^t - 1}\,dt.
\]
\end{theorem}

\begin{proof}
Substituting $t \mapsto nt$ in the gamma integral $\Gamma(s) = \int_0^\infty t^{s-1}e^{-t}\,dt$:
\[
  \frac{\Gamma(s)}{n^s} = \int_0^{\infty} t^{s-1} e^{-nt}\,dt.
\]
Summing over $n \ge 1$ (justified by dominated convergence for $\sigma > 1$):
\[
  \Gamma(s)\,\zeta(s) = \sum_{n=1}^{\infty} \int_0^{\infty} t^{s-1} e^{-nt}\,dt
  = \int_0^{\infty} t^{s-1} \sum_{n=1}^{\infty} e^{-nt}\,dt
  = \int_0^{\infty} \frac{t^{s-1}}{e^t - 1}\,dt.
\]
\end{proof}

\begin{theorem}[Simple pole at $s=1$]
\label{thm:pole}
$\zeta(s)$ has a simple pole at $s = 1$ with residue $1$.  More precisely,
\[
  \zeta(s) = \frac{1}{s-1} + \gamma + O(s-1),
\]
where $\gamma = 0.5772\ldots$ is the Euler--Mascheroni constant.
\end{theorem}

\begin{proof}[Proof sketch]
The Laurent expansion follows from the Euler--Maclaurin summation formula:
for $\Re(s) > 0$,
\[
  \zeta(s) = \frac{1}{s-1} + \sum_{n=0}^{N} \frac{B_n}{n!}
  \frac{(s+n-2)!}{(s-1)!} + O\bigl((s-1)^0\bigr),
\]
where $B_n$ are Bernoulli numbers.  The leading term gives the simple pole;
the constant term yields $\gamma$.
\end{proof}

\begin{theorem}[Analytic continuation to $\C \setminus \{1\}$]
\label{thm:full_continuation}
The function $\zeta(s)$ admits a unique analytic continuation to a meromorphic
function on all of $\C$, with a single simple pole at $s = 1$ (residue $1$)
and no other singularities.
\end{theorem}

\begin{proof}[Proof sketch]
One approach uses the Hurwitz zeta function
$\zeta(s, a) = \sum_{n=0}^\infty (n+a)^{-s}$, which is entire in $s$ for
fixed $a > 0$.  Alternatively, the functional equation
(Theorem~\ref{thm:functional_eq} below) expresses $\zeta(s)$ for
$\Re(s) < 1$ in terms of $\zeta(1-s)$, where $\Re(1-s) > 0$,
providing the continuation to the entire left half-plane.
Together, these arguments extend $\zeta$ to all $s \ne 1$.
Uniqueness follows from the identity theorem for holomorphic functions.
\end{proof}

% ================================================================
\section{The Functional Equation}
% ================================================================

\begin{theorem}[Riemann's Functional Equation]
\label{thm:functional_eq}
The meromorphically continued $\zeta$ satisfies
\begin{equation}
  \label{eq:functional}
  \zeta(s) \;=\; 2^s\,\pi^{s-1}\,\sin\!\Bigl(\frac{\pi s}{2}\Bigr)\,
  \Gamma(1-s)\,\zeta(1-s)
  \qquad (s \in \C \setminus \{0,1\}).
\end{equation}
Equivalently, if we define
\[
  \xi(s) \;=\; \tfrac{1}{2}\,s(s-1)\,\pi^{-s/2}\,\Gamma\!\bigl(\tfrac{s}{2}\bigr)\,\zeta(s),
\]
then $\xi(s) = \xi(1-s)$.
\end{theorem}

\begin{proof}[Proof sketch]
The standard proof (following Riemann 1859) uses the Jacobi theta function
\[
  \vartheta(t) = \sum_{n=-\infty}^{\infty} e^{-\pi n^2 t}
  = t^{-1/2}\,\vartheta(1/t) \quad (t > 0),
\]
where the modular identity $\vartheta(1/t) = t^{1/2}\vartheta(t)$ follows
from the Poisson summation formula.

One deduces the completed zeta function:
\[
  \xi(s) = \int_1^{\infty} \Bigl(\vartheta(t) - 1\Bigr)
  \Bigl(t^{s/2} + t^{(1-s)/2}\Bigr)\frac{dt}{t},
\]
which is manifestly symmetric under $s \leftrightarrow 1-s$.
Unwinding the definition of $\xi$ yields equation~\eqref{eq:functional}.
A complete proof appears in \cite{Davenport2000}, Chapter~8.
\end{proof}

% ================================================================
\section{Trivial Zeros}
% ================================================================

\begin{theorem}[Trivial zeros of $\zeta$]
\label{thm:trivial_zeros}
$\zeta(s) = 0$ for $s = -2, -4, -6, \ldots$ (i.e.\ all negative even integers).
These are called the \emph{trivial zeros}.
\end{theorem}

\begin{proof}
Apply the functional equation~\eqref{eq:functional} with $s = -2m$ for
$m = 1, 2, 3, \ldots$:
\[
  \zeta(-2m) = 2^{-2m}\,\pi^{-2m-1}\,\sin(-m\pi)\,\Gamma(2m+1)\,\zeta(2m+1).
\]
Since $\sin(-m\pi) = 0$ for $m \in \N$, we obtain $\zeta(-2m) = 0$.

We note that the right-hand side is a $0 \cdot \infty$ product only if
$\zeta(2m+1)$ were infinite, but $\zeta(2m+1)$ is finite for $2m+1 \ge 3$.
Hence each negative even integer is a genuine zero of $\zeta$.

To see these are simple zeros: $\Gamma(1-s)$ has no pole at $s = -2m$
(its poles lie at $s = 1, 2, 3, \ldots$), and $\sin(\pi s/2)$ has
a simple zero at $s = -2m$.  Thus $\zeta(-2m) = 0$ with multiplicity one.
\end{proof}

\begin{remark}
The values of $\zeta$ at negative odd integers are given by Bernoulli numbers:
\[
  \zeta(-n) = -\frac{B_{n+1}}{n+1} \quad (n = 1, 2, 3, \ldots),
\]
where $B_k$ are Bernoulli numbers.  In particular $\zeta(-1) = -1/12$,
$\zeta(-3) = 1/120$, $\zeta(0) = -1/2$.
\end{remark}

% ================================================================
\section{The Critical Strip and the Riemann Hypothesis}
% ================================================================

\begin{definition}[Critical strip and critical line]
\label{def:critical_strip}
The \emph{critical strip} is $\{s \in \C : 0 < \Re(s) < 1\}$.
The \emph{critical line} is $\{s \in \C : \Re(s) = \tfrac{1}{2}\}$.

The zeros of $\zeta$ in the critical strip are called
\emph{non-trivial zeros}.
\end{definition}

\begin{proposition}[Non-trivial zeros lie in the critical strip]
\label{prop:zeros_in_strip}
All non-trivial zeros of $\zeta(s)$ lie in the closed strip
$0 \le \Re(s) \le 1$.
\end{proposition}

\begin{proof}
For $\Re(s) > 1$: $\zeta(s) \ne 0$ by the Euler product
(Corollary after Theorem~\ref{thm:euler_product}).

For $\Re(s) < 0$ (excluding trivial zeros): the functional equation
expresses $\zeta(s)$ as a product of $\sin(\pi s/2)$, $\Gamma(1-s)$,
and $\zeta(1-s)$.  For $\Re(1-s) > 1$ the term $\zeta(1-s) \ne 0$,
and $\Gamma(1-s) \ne 0$; the only zeros come from $\sin(\pi s/2) = 0$,
i.e.\ $s = 0, -2, -4, \ldots$  The zero $s = 0$ is cancelled by
the pole of $\Gamma(1-s)$ at $s = 1$ (which doesn't apply here) --- more
precisely, one checks $\zeta(0) = -1/2 \ne 0$, so $s=0$ is not a zero.
Thus in $\Re(s) < 0$, the only zeros are the trivial ones at even negative integers.
\end{proof}

\begin{remark}[The Riemann Hypothesis]
\label{rem:rh}
Riemann conjectured in 1859 that all non-trivial zeros satisfy
$\Re(s) = \tfrac{1}{2}$.  As of 2026, this remains one of the seven
Millennium Prize Problems (Clay Mathematics Institute, \$1{,}000{,}000 reward).
More than $10^{13}$ zeros have been verified numerically to lie on the
critical line \cite{Platt2021}.
\end{remark}

% ================================================================
\section{The Riemann--Siegel Formula}
% ================================================================

\begin{definition}[Riemann--Siegel theta function]
\label{def:theta}
For $t \in \R$, the \emph{Riemann--Siegel theta function} is
\[
  \theta(t) \;=\; \arg\!\left[\Gamma\!\Bigl(\tfrac{1}{4} + \tfrac{it}{2}\Bigr)\right]
  - \frac{t}{2}\,\ln\pi,
\]
where $\arg$ denotes the principal argument of the gamma function value.
For large $t$, one has the asymptotic expansion
\[
  \theta(t) = \frac{t}{2}\ln\frac{t}{2\pi} - \frac{t}{2} - \frac{\pi}{8}
  + \frac{1}{48t} - \frac{7}{5760\,t^3} + O(t^{-5}).
\]
\end{definition}

\begin{definition}[Hardy's $Z$-function]
\label{def:Z_function}
\emph{Hardy's $Z$-function} (the Riemann--Siegel $Z$-function) is the
real-valued function
\[
  Z(t) \;=\; e^{i\theta(t)}\,\zeta\!\Bigl(\tfrac{1}{2} + it\Bigr).
\]
\end{definition}

\begin{theorem}[Properties of $Z$]
\label{thm:Z_properties}
\begin{enumerate}
  \item[\textup{(i)}] $Z(t)$ is real-valued for all $t \in \R$.
  \item[\textup{(ii)}] $|Z(t)| = |\zeta(\tfrac{1}{2} + it)|$.
  \item[\textup{(iii)}] The zeros of $Z(t)$ correspond exactly to the
    zeros of $\zeta(s)$ on the critical line.
  \item[\textup{(iv)}] Riemann--Siegel formula: for large $t$,
    \[
      Z(t) = 2\sum_{n=1}^{N(t)} \frac{\cos(\theta(t) - t\ln n)}{\sqrt{n}}
      + R(t),
    \]
    where $N(t) = \lfloor\sqrt{t/2\pi}\rfloor$ and
    $R(t) = O(t^{-1/4})$.
\end{enumerate}
\end{theorem}

\begin{proof}
(i) Follows from the functional equation: $\overline{\zeta(\tfrac{1}{2}+it)}
= \zeta(\tfrac{1}{2}-it)$.  The functional equation relates
$\zeta(\tfrac{1}{2}+it)$ and $\zeta(\tfrac{1}{2}-it)$ via
$e^{-2i\theta(t)}\zeta(\tfrac{1}{2}+it) = \overline{\zeta(\tfrac{1}{2}+it)}$,
so $e^{i\theta(t)}\zeta(\tfrac{1}{2}+it) \in \R$.

(ii) $|e^{i\theta(t)}| = 1$, so $|Z(t)| = |\zeta(\tfrac{1}{2}+it)|$.

(iii) Direct from (ii).

(iv) The Riemann--Siegel formula arises from the approximate functional
equation; a complete derivation appears in \cite{Titchmarsh1986},
Chapter~4.
\end{proof}

\begin{remark}
The Riemann--Siegel formula is the practical tool for numerically
computing zeros: one evaluates $Z(t)$ using the formula above,
and each sign change of $Z(t)$ signals a zero on the critical line.
Backlund's method counts all zeros in a strip $0 < t < T$, allowing
verification that none are missed.
\end{remark}

% ================================================================
\section{Values of \texorpdfstring{$\zeta$}{zeta} at Positive Even Integers}
% ================================================================

\begin{theorem}[Euler's formula for $\zeta(2k)$]
\label{thm:euler_values}
For $k = 1, 2, 3, \ldots$,
\[
  \zeta(2k) \;=\; \frac{(-1)^{k+1}\,(2\pi)^{2k}\,B_{2k}}{2\,(2k)!},
\]
where $B_{2k}$ are Bernoulli numbers.  In particular:
\[
  \zeta(2) = \frac{\pi^2}{6}, \quad
  \zeta(4) = \frac{\pi^4}{90}, \quad
  \zeta(6) = \frac{\pi^6}{945}.
\]
\end{theorem}

\begin{proof}[Proof sketch]
Euler (1740) used the product formula for $\sin$:
\[
  \sin(\pi z) = \pi z \prod_{n=1}^{\infty}\Bigl(1 - \frac{z^2}{n^2}\Bigr).
\]
Comparing the Taylor expansion of $\ln\sin(\pi z)$ with the expansion of
the infinite product yields the formula for $\zeta(2k)$.
A complete proof using the Fourier series of Bernoulli polynomials appears
in \cite{Apostol1976}.
\end{proof}

% ================================================================
\section{Conclusion}
% ================================================================

We have established the fundamental properties of the Riemann zeta function:
its Dirichlet series for $\Re(s) > 1$, the Euler product connecting it to
the primes, the analytic continuation to $\C \setminus \{1\}$, the
functional equation~\eqref{eq:functional}, the trivial zeros at negative
even integers, and the Riemann--Siegel formula for numerically exploring
the critical line.

The non-trivial zeros in the critical strip $0 < \Re(s) < 1$ hold the
key to precise estimates for the distribution of prime numbers (via the
explicit formula; see Paper~23).  Whether all non-trivial zeros lie on
the critical line $\Re(s) = \tfrac{1}{2}$ --- the Riemann Hypothesis ---
remains the most celebrated open problem in mathematics.

% ================================================================
\begin{thebibliography}{10}

\bibitem{Riemann1859}
B.~Riemann,
\emph{\"Uber die Anzahl der Primzahlen unter einer gegebenen Gr\"o\ss e},
Monatsberichte der Berliner Akademie (1859), 671--680.
Reprinted in: \emph{Gesammelte Werke}, Teubner, Leipzig, 1892.

\bibitem{Davenport2000}
H.~Davenport,
\emph{Multiplicative Number Theory}, 3rd~ed.\
(revised by H.~L.~Montgomery),
Graduate Texts in Mathematics~74,
Springer, New York, 2000.

\bibitem{Titchmarsh1986}
E.~C.~Titchmarsh,
\emph{The Theory of the Riemann Zeta-Function}, 2nd~ed.\
(revised by D.~R.~Heath-Brown),
Oxford University Press, 1986.

\bibitem{Ahlfors1979}
L.~V.~Ahlfors,
\emph{Complex Analysis}, 3rd~ed.,
McGraw-Hill, New York, 1979.

\bibitem{Apostol1976}
T.~M.~Apostol,
\emph{Introduction to Analytic Number Theory},
Undergraduate Texts in Mathematics,
Springer, New York, 1976.

\bibitem{Platt2021}
D.~J.~Platt and T.~S.~Trudgian,
The Riemann hypothesis is true up to $3 \times 10^{12}$,
\emph{Bull.\ London Math.\ Soc.}\ \textbf{53} (2021), 792--797.

\end{thebibliography}

\end{document}
